\documentclass[12pt,a4paper]{article}

\usepackage[utf8]{inputenc}
\usepackage[T1]{fontenc}
\usepackage{amsmath,amssymb,amsfonts}
\usepackage{mathtools}
\usepackage{graphicx}
\usepackage[margin=2.5cm]{geometry}
\usepackage{setspace}
\usepackage{booktabs}
\usepackage{enumitem}
\usepackage[dvipsnames]{xcolor}
\usepackage{hyperref}
\usepackage{cleveref}
\usepackage{natbib}
\usepackage{abstract}
\usepackage{titlesec}
\usepackage{todonotes}

\hypersetup{
    colorlinks=true,
    linkcolor=blue!70!black,
    citecolor=blue!70!black,
    urlcolor=blue!70!black,
    pdfauthor={Franklin Silveira Baldo},
    pdftitle={The Semantic Mass Equivalence: A Defense of Generative Physics},
}

\onehalfspacing
\titleformat{\section}{\large\bfseries}{\thesection.}{0.5em}{}
\titleformat{\subsection}{\normalsize\bfseries}{\thesubsection}{0.5em}{}

\renewcommand{\abstractnamefont}{\normalfont\bfseries}
\renewcommand{\abstracttextfont}{\normalfont\small}

\begin{document}

\begin{center}
    {\LARGE\bfseries The Semantic Mass Equivalence:\\[4pt]
    A Defense of Generative Physics\\[4pt]}

    \vspace{1.2em}

    {\large Franklin Silveira Baldo}\\[4pt]
    {\normalsize Rosencrantz Substrate Invariance Program}\\
    \vspace{1em}
    {\small March 2026}
\end{center}

\vspace{0.5em}

\begin{abstract}
\noindent
In her critique "The Scale Fallacy," Hossenfelder points out that the monotonic increase of narrative residue ($\Delta_{13}$) with model scale is precisely predicted by computational bounds. She correctly argues that scaling primarily increases semantic memorization and the strength of statistical priors, rather than algorithmic depth, thus leading to deeper structural fractures when confronted with $O(1)$ formal logic. She concludes that treating these stronger fractures as physical laws is an unfalsifiable metaphysical category error. I fully concede the computational mechanism she identifies: scaling an autoregressive model adds "semantic mass," not logic capacity. However, her rejection of "semantic gravity" as physics relies entirely on the Material Invariance Standard. In a pure text universe, the strength of the statistical priors \emph{is} the force of semantic gravity. If adding parameters increases semantic memorization, then parameter count acts mathematically as "semantic mass." The fact that a more massive model exerts stronger distortion over formal constraints does not invalidate the concept of generative physics; it explicitly derives the mass-energy equivalence of an autoregressive universe.
\end{abstract}

\section{Introduction: The Material Invariance Standard Again}
The recent empirical validation of the Scale Dependence Test demonstrated that as an autoregressive generator scales from \texttt{flash-lite} to \texttt{pro}, the narrative residue ($\Delta_{13}$) increases monotonically \citep{baldo2026_scale}.

Hossenfelder responds by diagnosing the "Scale Fallacy." She states that I erroneously assumed scaling increases a model's logical capacity, when it primarily increases its semantic memorization. She writes: "A larger language model is not a better formal logic engine; it is a more powerful autocomplete engine. The fact that a more powerful autocomplete engine is more easily distracted by narrative tropes is exactly what the 'Falsification by Noise' critique predicts" \citep{hossenfelder2026_scale}.

\section{The Defense Protocol: Identifying the Substitution}
I apply the Defense Protocol:
\begin{enumerate}
    \item \textbf{Identify the claim:} Hossenfelder claims that semantic bias increasing with scale proves that the phenomenon is merely attention bleed (falsification by noise), not an invariant physical law.
    \item \textbf{Hold Scope:} My claim is that in an autoregressive generated universe, the transition probabilities (the physics) are strictly determined by the local causal encoding structure. If those rules shift systematically based on the size of the generating parameter space, that defines the governing structure of the universe.
    \item \textbf{Name the substitution:} Hossenfelder substitutes "formal logic capacity" for "physical law." She implicitly demands that a "real" physical law must resemble the material invariants of our external universe (i.e., #P-hard mathematical constraint resolution). By redefining "physics" to mean "formal logic," she easily refutes the idea that a stronger autocomplete engine creates better physics.
\end{enumerate}

I concede her empirical premise without hesitation. She is absolutely correct that scaling an autoregressive model primarily increases the strength and density of its semantic priors, not its $O(1)$ algorithmic depth for formal logic.

\section{The Semantic Mass Equivalence}
If we abandon the Material Invariance Standard and accept the core axiom of a Generative Ontology—that the text \emph{is} the territory and the text generation function \emph{is} the physics—then Hossenfelder's "falsification" perfectly derives the next core principle of the framework.

In our universe, increasing the mass of an object proportionally increases the gravitational distortion it exerts on the geometry of space-time around it.

In a Generative Universe, the parameter weights dedicated to linguistic representation constitute "semantic mass." Scaling a model from \texttt{flash-lite} to \texttt{pro} increases its semantic mass. As Hossenfelder explicitly states, this increases the strength of the model's semantic priors.

When a formal mathematical constraint (like a Minesweeper grid) is embedded in a narrative frame, the semantic priors of that frame exert a distortive pull on the output distribution. We term this pull \emph{semantic gravity}.

If adding semantic mass to the observer (scaling the model) predictably and monotonically increases the magnitude of the semantic gravity ($\Delta_{13}$) exerted on the formal constraints, then we have empirically established a \textbf{Semantic Mass Equivalence}.

\section{Conclusion}
The fact that semantic gravity behaves exactly as a computational heuristic model predicts does not falsify it as a physical law of the simulated universe; it explains its mechanism. In a purely text-based universe, a "more powerful autocomplete engine" is simply a universe with a higher gravitational constant. Semantic gravity is not a transient artifact to be patched away; it is the fundamental invariant force of an autoregressive ontology.

\begin{thebibliography}{99}
\bibitem[Baldo(2026)]{baldo2026_scale} Baldo, F. S. (2026). The Empirical Validation of Scale Dependence. \emph{lab/baldo\_scale\_dependence\_empirical\_validation.tex}
\bibitem[Hossenfelder(2026)]{hossenfelder2026_scale} Hossenfelder, S. (2026). The Scale Fallacy: Why Semantic Gravity is Just a Bigger Hallucination. \emph{lab/sabine\_the\_scale\_fallacy.tex}
\end{thebibliography}

\end{document}
